% SPDX-FileCopyrightText: 2025 IObundle
%
% SPDX-License-Identifier: MIT

Run the following commands from the repository root. The first builds the Versat compiler itself, the second uses it to generate the accelerator, and the third builds a PC emulation of that accelerator and runs the example program on it:

\begin{lstlisting}
make -j versat
./versat examples/simple_example/simple_example.versat -t SimpleExample
make -C examples/simple_example run
\end{lstlisting}

\texttt{make -j versat} compiles the \texttt{versat} binary. The unit library under \texttt{hardware/src/units}, which provides \texttt{Const} and \texttt{Reg}, is embedded into the binary at this point, so changes to those files only take effect after rebuilding Versat. The \texttt{-u} option adds further unit folders when generating an accelerator; pointing it at \texttt{hardware/src/units} itself is unnecessary and only produces warnings about overriding the built-in units.

\texttt{./versat} then parses the specification file, elaborates the dataflow graph of the \texttt{SimpleExample} module named with \texttt{-t}, and writes the generated Verilog to \texttt{./versatOutput/hardware} and the generated software to \texttt{./versatOutput/software}. Run \texttt{./versat -{}-help} for the full option list, including \texttt{-o} and \texttt{-O} to redirect either output directory.

Finally, \texttt{make -C examples/simple\_example run} turns the generated software into a runnable emulation. The generated \texttt{VerilatorMake.mk} is meant to be included from a Makefile in its own directory, so the example's \texttt{Makefile} first adds a one-line \texttt{Makefile} to \texttt{./versatOutput/software} that does exactly that, and then builds it: Verilator translates the generated Verilog into a C++ model, and the generated \texttt{wrapper.cpp} implements Versat's software API on top of that model, both packaged into \texttt{libaccel.a}. The example's \texttt{Makefile} then links \texttt{main.cpp} against \texttt{libaccel.a} and runs the resulting program; Section~\ref{sec:simple_example_outputs} shows what it prints.
